\documentclass[a4paper,10pt]{article}
\usepackage{cw2026-report}

\newcommand{\ReportTitle}{Sensor footprint and channel matching}
\newcommand{\ReportSubtitle}{What the reflectance array can resolve, and what the twin must model}

\newcommand{\StudentID}{[7-digit student ID]}

\begin{document}
\ReportHeader

\begin{multicols*}{2}
\raggedcolumns

\section*{Subsystem and investigation aim}
Our 3Pi+ scans a printed sheet row by row and rebuilds a picture from the five
downward-facing reflectance sensors. Each reading is the time in microseconds for
a capacitor to discharge through a phototransistor, so a bright surface returns a
short time and dark ink a long one. New snapshots arrive every 10\,ms.

This report characterises the array's spatial footprint and its channel-to-channel
matching, because both are missing from the digital twin and neither can be
inferred from the controller. The question is what spatial and tonal resolution
the real array delivers at the settings we scan with, and what the twin must
represent to match it.

\section*{Measurement approach}
Three measurements were taken. The robot was stationary and the surface was the
A3 test sheet unless stated.

\textbf{Geometry.} Each channel's position was read from a photograph of the
underside against a rule, relative to the wheel-axle centre. The measured offsets
are 28 to 41\,mm ahead and within 27\,mm of the centre line, for DN1 to DN5, with
$\pm$2\,mm uncertainty from pixel resolution and parallax.

\textbf{Response.} The array was placed over plain white paper and over solid
black ink, and the per-channel times were read from the 10\,ms snapshot stream.
White gives about 400\,$\mu$s and black about 2000\,$\mu$s, the direction expected
from brighter reflection. Two prints bound the usable range but cannot establish
linearity, and they say nothing about saturation, so the midpoint, 1200\,$\mu$s, is
reported as the midpoint of the range and not as a calibrated 50 percent tone.

\textbf{Footprint.} The footprint is the weakest of the three measurements. Each
reading averages the paper under the channel, so a printed edge comes back
softened. We chose the Gaussian width $\sigma$ that, when used to deblur our scans,
removed the softening without ringing, and halos appeared above about 4\,mm. That
gives $\sigma\approx3$\,mm with a plausible range of 2 to 4\,mm. It is an estimate
from the artefact it explains, not an independent measurement.

\ReportPlot{fig_sensor.pdf}{Measured sensor geometry and response}
{Left, measured channel positions relative to the axle centre, $\pm$2\,mm, with
the estimated 3\,mm spot drawn to scale and the 13\,mm fore-aft stagger marked.
Right, measured discharge time against printed brightness. The two prints bound
the range and do not establish a linear response.}
{fig:sensor}

\section*{Results and interpretation}
The array spans 54\,mm laterally and 13\,mm fore and aft. The stagger matters.
A channel 13\,mm ahead of another samples a different part of the sheet, so the
five channels interlace rather than share a row, and a rebuild that treats them as
co-located shifts every feature by 13\,mm, which is nearly twice the 7\,mm row
pitch we scan with.

At $\sigma=3$\,mm the width from 10 to 90 percent of an edge is
$2.56\times3\approx7.7$\,mm. That is 1.1 times our row pitch and 27 times the
0.28\,mm encoder resolution. The honest reading is that the spot and the row
spacing are comparable, so vertical resolution is set by both together and this
measurement alone cannot say which dominates. It does explain why a printed
triangle comes back as a soft mound, since its corners are comparable to the spot.
At $\sigma=2$\,mm the width would be 5.1\,mm and at 4\,mm, 10.2\,mm, so the
estimate of $\sigma$ dominates the uncertainty.

The tonal result is weaker and is reported as such. Between 400 and 2000\,$\mu$s
the two prints give a straight line, which is not evidence of a linear sensor. Our
rebuilds with a single shared brightness range show stripes on uniform paper, which
is direct evidence that the channels differ in offset and gain. We have not
quantified those differences, so the model's assumption of five identical channels
is known to be wrong but not yet replaced by measured values.

\section*{Implications for the digital twin}
The twin treats each channel as a point sensor with identical levels. Two changes
follow from the evidence. Each channel should be a Gaussian spot of width $\sigma$
at its measured offset, with its own white and black times. The offset part is
supported by the stripes. The $\sigma$ part is supported only by the deblur
estimate, so it should be implemented as a parameter with a stated range rather
than a fixed number.

The consequence of leaving them out is systematic, not random. A point-sensor twin
produces edges sharper than the robot can and rewards a finer row pitch than the
array can use, so scan settings chosen in simulation will underperform at the
bench.

A suitable check needs an observation not used for the fit. Scan one straight
black and white edge at three speeds and measure the width from 10 to 90 percent.
The footprint is fixed while any delay grows with speed, so a constant width
confirms the 3\,mm $\sigma$ and a growing width shows that delay dominates. A
printed step wedge, measured with repeats, would then give a gain and an offset per
channel.

\EvidenceNote{Measurements and processing are in
\texttt{ModelParameters.pde}, which holds the channel offsets,
\texttt{sensorFootprintSigmaMm} and the white and black times, and in the scan
logs exported by the dashboard, which carry \texttt{dn1} to \texttt{dn5} at
100\,Hz. The deblur fit and the sensor-position table are in the dashboard Image
settings panel. The knife-edge and step-wedge tests have not been run.
References: [1] G.~D. Boreman, \emph{Modulation Transfer Function}, SPIE Press,
2001, ch.~3. [2] Pololu, \emph{QTR RC-type sensor application note},
pololu.com/docs/0J13.}

\end{multicols*}
\end{document}
